\documentclass[12pt]{article}

% --------------------
% Paquetes
% --------------------
\usepackage[spanish]{babel}
\usepackage[utf8]{inputenc}
\usepackage[T1]{fontenc}
\usepackage{setspace}
\usepackage{geometry}
\usepackage{times}
\usepackage{graphicx}
\usepackage{amsmath}
\usepackage{soul}
\usepackage{tabularx}
\usepackage{booktabs}

% --------------------
% Márgenes
% --------------------
\geometry{
  letterpaper,
  left=2.5cm,
  right=2.5cm,
  top=2.5cm,
  bottom=2.5cm
}

% --------------------
% Documento
% --------------------
\begin{document}
\thispagestyle{empty}

% --------------------
% Portada
% --------------------
\begin{center}

{\LARGE \textbf{Instituto Politécnico Nacional}}\\[0.5cm]
{\large Escuela Superior de Cómputo}\\[1.5cm]



{\Huge \textbf{Procesamiento de Lenguaje Natural}}\\[0.5cm]
{\large Segundo Parcial}\\[2cm]

{\LARGE \textbf{Tareas}}\\[0.5cm]

\end{center}

\vfill

\begin{minipage}{0.55\textwidth}
\textbf{Alumno:}\\
Sánchez García Miguel Alexander\\[1cm]

\textbf{Grupo:} 6AV1\\
\textbf{Carrera:} Licenciatura en Ciencia de Datos
\end{minipage}
\hfill
\begin{minipage}{0.35\textwidth}
\textbf{Profesor:}\\
Ortiz Castillo Marco Antonio\\[1cm]

\textbf{Fecha:}\\
9 de enero, 2026
\end{minipage}

\newpage

% --------------------
% Contenido
% --------------------
\setstretch{1.3}


\section*{\hl{One-Hot Encoding}}

Dado el vector de palabras, se deben eliminar y contar aquellas palabras que estén repetidas.

\subsection*{Ejemplo}

Sea el siguiente vector de palabras:

\begin{equation}
\text{vector} = [\text{``Marco''}, \text{``Marco''}, \text{``Antonio''}, \text{``Ortiz''}, \text{``Ortiz''}]
\end{equation}

El primer paso consiste en crear un vector que contenga únicamente las palabras únicas.

\subsection*{Vector de palabras únicas}

\begin{equation}
\text{vector}_\text{unic} = [\text{``Marco''}, \text{``Antonio''}, \text{``Ortiz''}]
\end{equation}

\section*{\hl{Bolsa de Palabras (Bag of Words, BoW)}}

La Bolsa de Palabras es una técnica de cuantificación de las palabras.  
Junto con TF--IDF, permite dotar al documento de valores numéricos; también se dice que extrae características del documento.

Es un paso previo a otros algoritmos, y se da por hecho que ya ha existido un pre-procesamiento del texto.  
Se puede expresar mediante una matriz y es un pre-algoritmo útil para la obtención de palabras clave.

\subsection*{Ejemplo}

Sea el siguiente documento:

\begin{quote}
Tenemos una invitada que habla francés.  
El francés no toma como invitada a cualquier persona, debe ser de habla francesa, puesto que cualquier otra persona será excluida.
\end{quote}

Por cada palabra se crea un vector de dimensión $n$, donde $n$ es el total de palabras únicas del documento.

\subsection*{Vector de palabras únicas}

\begin{equation}
\begin{aligned}
\text{vector}_{\text{único}} = [&\text{``Tenemos''}, \text{``invitada''}, \text{``habla''}, \text{``francés''}, \text{``no''}, \text{``toma''}, \\
&\text{``cualquier''}, \text{``persona''}, \text{``debe''}, \text{``ser''}, \text{``francesa''}, \\
&\text{``puesto''}, \text{``otra''}, \text{``será''}, \text{``excluida''}]
\end{aligned}
\end{equation}

\begin{equation}
n = 15
\end{equation}

\subsection*{Representación vectorial}

Cada palabra se representa como un vector donde la posición correspondiente indica el número de apariciones de dicha palabra en el documento.

\begin{align}
\text{Tenemos} &= [1, 0, 0, 0, \ldots, 0] \\
\text{invitada} &= [0, 2, 0, 0, \ldots, 0] \\
\text{habla} &= [0, 0, 1, 0, \ldots, 0] \\
\vdots &
\end{align}

\subsection*{Observaciones}

Para reducir la dimensionalidad del modelo, se recomienda utilizar algoritmos de reducción de dimensión como PCA o SVD.

Asimismo, también puede emplearse lematización como parte del pre-procesamiento.

\section*{\hl{Bolsa de Palabras en múltiples documentos}}

La bolsa de palabras suele usarse en varios documentos, en donde se obtendrá un solo vector de palabras que caracterice a todo el conjunto de documentos.  
Este vector se obtiene a partir de un vector único de palabras construido considerando todos los documentos.

\subsection*{Ejemplo}

Se tienen los siguientes documentos:

\begin{itemize}
    \item \textbf{Doc 1:} José piensa que canta bello.
    \item \textbf{Doc 2:} José podría ir a clases de canto.
    \item \textbf{Doc 3:} Su amigo podría decirle la verdad de cómo canta.
\end{itemize}

\subsection*{Matriz BoW}

\paragraph{a) Vector único de palabras}

A partir de los documentos anteriores, se construye el siguiente vector único de palabras:

\begin{equation}
\begin{aligned}
\text{vector}_{\text{único}} = [&\text{``José''}, \text{``piensa''}, \text{``canta''}, \text{``bello''}, \text{``podría''}, \text{``ir''}, \\
&\text{``clases''}, \text{``canto''}, \text{``su''}, \text{``amigo''}, \text{``decirle''}, \\
&\text{``verdad''}, \text{``cómo''}]
\end{aligned}
\end{equation}

\paragraph{b) Vectores por documento}

Cada documento se representa como un vector donde cada posición indica el número de veces que aparece la palabra correspondiente del vector único.

\begin{align}
\vec{v}_{\text{doc}_1} &= [1, 1, 1, 1, 0, 0, 0, 0, 0, 0, 0, 0, 0] \\
\vec{v}_{\text{doc}_2} &= [1, 0, 0, 0, 1, 1, 1, 2, 0, 0, 0, 0, 0] \\
\vec{v}_{\text{doc}_3} &= [0, 0, 1, 0, 1, 0, 0, 0, 1, 1, 1, 1, 1]
\end{align}

\section*{\hl{Embeddings}}

Son representaciones numéricas puestas en un vector de dimensión reducida que captura la relación semántica y sintáctica de las palabras. A diferencia de otros métodos que la representación es binaria, generalmente usan valores entre $-1$ y $1$ con todos los números reales que puede haber entre estos. 

\textbf{marco} = $[-0.35, .55, .777]$

\vspace{1em}

\section*{\hl{Similitud Coseno}}

Es una métrica de similitud que, en lugar de \hl{evaluar la longitud} o tamaño, evalúa la dirección en la que apuntan dos vectores.

Especialmente útil cuando se tienen representaciones vectoriales de alta dimensión.

\vspace{1em}

\begin{minipage}[t]{0.5\textwidth}
    La similitud coseno (Que en realidad mide el ángulo entre dos vectores) regresa un valor entre $-1$ y $1$.
    
    \vspace{1em}
    Cuando el valor es $1$ se dice que las palabras son totalmente similares, cuando es $0$, los vectores son ortogonales y por tanto diferentes. Finalmente cuando es $-1$, indica que sus vectores son similares pero que apuntan a direcciones contrarias.
    
    \vspace{1em}
    \textbf{Ejemplo:} Perro - Gato
\end{minipage}
\hfill
\begin{minipage}[t]{0.4\textwidth}
    \centering
    \[ \cos(A, B) = \frac{A \cdot B}{\|A\| \cdot \|B\|} \]
    
    \flushleft
    \textbf{Ventajas:}
    \begin{itemize}
        \item Recomendado para espacios de alta dimensión.
        \item Insensibles a la magnitud.
        \item Aplicable en todo dominio.
    \end{itemize}
\end{minipage}


\section*{\hl{Expresiones regulares}}

Son un conjunto de caracteres que definen una regla de coincidencia para obtener un conjunto específico de palabras, por ejemplo, fechas, verbos, números telefónicos, etc...

\vspace{1em}

Permiten eliminar signos de puntuación, números o símbolos:
\begin{itemize}
    \item Tokenizar por espacios o puntos.
    \item Extraer información.
    \item Filtrar palabras.
\end{itemize}

\vspace{2em}

\begin{center}
\begin{tabular}{l @{\hspace{4em}} l}
    \textbf{Expresión} & \textbf{Significado} \\
    \rule{0pt}{4ex} % Espacio extra
    \verb|\d+| & 1 o más cifras \\
    \rule{0pt}{3ex}
    \verb|\b\w+\b| & una palabra completa \\
    \rule{0pt}{3ex}
    \verb|[A-Z][a-z]+| & palabra inicia en mayúscula \\
    \rule{0pt}{3ex}
    \verb|https: //\S+| & URL \\
    \rule{0pt}{3ex}
    \verb|@\w+| & menciones en redes \\
\end{tabular}
\end{center}

\vspace{2em}

Librerías utilizadas para hacer búsqueda por expresiones regulares es ``re'', NLTK y SpaCy.


\section*{\centering \hl{Tarea: Diferencia entre TSNE y PCA}}

Tanto PCA (Análisis de componentes principales) como t-SNE (Incrustación estocástica de vecinos distribuida en t) son utilizados para reducción de dimensionalidad, pero tienen objetivos distintos.

\vspace{1.5em}

\begin{tabularx}{\textwidth}{p{3cm} X X}
    \toprule
    & \textbf{PCA} & \textbf{t-SNE} \\
    \midrule
    \textbf{Tipo de método} & Lineal & No-Lineal \\
    \addlinespace
    \textbf{Objetivo principal} & Maximiza varianza & Preserva vecindarios \\
    \addlinespace
    \textbf{Qué se interpreta} & Componentes (Qué variables pesan) & Ejes no interpretables \\
    \addlinespace
    \textbf{Uso típico} & Feature engineering, evita colinealidad, acelerar modelos, análisis exploratorio & Visualización de 2D/3D datos de alta dimensión para ver agrupamientos locales. \\
    \addlinespace
    \textbf{Determinismo} & Determinista (dado el preprocesamiento) & Estocástico (depende de semilla e hiperparámetros) \\
    \addlinespace
    \textbf{Hiperparámetros} & \# componentes & Perplexity, learning rate, iteraciones, semilla \\
    \addlinespace
    \textbf{Escalabilidad} & Muy eficiente, escala bien & Más costoso y puede ser lento en datasets grandes \\
    \bottomrule
\end{tabularx}


\end{document}